\subsection{Linear Programming Definition}

The general form of a linear program is given 
in Equation \ref{eq:genlp}.

\begin{align}
    &\text{min}_{x} / \text{max}_x &\zeta = c_1x_1 + c_2x_2 + \dots + c_nx_n + c_0 \\
    &\text{s.t.} &b_1^{\text{L}} \leq a_{1,1} x_1 + a_{1,2}x_2 + \dots + a_{1,n}x_n \leq b_1^{\text{U}} \\
    &&\vdots \\
    &&b_m^{\text{L}} \leq a_{m,1}x_1 + a_{m,2}x_2 + \dots + a_{m,n}x_n \leq b_m^{\text{U}} \\
    &&l_1 \leq x_1 \leq u_1 \\
    &&\vdots \\
    &&l_n \leq x_n \leq u_n.
\end{align}
\label{eq:genlp}

Definitions for solutions:
\begin{itemize}
    \item The variables $x_1, x_2, \dots, x_n$
        are called \emph{decision variables},
        and $a$, $b^{\text{L}}$, $b^{\text{U}}$,
        $c$, $l$, and $u$ are given \emph{parameters}.
    \item A solution $x = (x_1, x_2, \dots, x_n)$ is a \emph{feasible solution} if it satisfies all
            constraints. A solution is an \emph{infeasible solution} if it violates
            any constraint.
    \item The set of all feasible solutions is called the \emph{feasible region} or
        \emph{feasible set}.
    \item $l_i$ and $u_i$ are called the lower bound and upper bound on variable
        $x_i$, respectively.
    \item If $l_i$ = $-\infty$ and $u_i = +\infty$, 
        then $x_i$ is a free or unrestricted variable.
\end{itemize}

Definitions for constraints:
\begin{itemize}
    \item The symbol 's. t.' is short for 'subject to' 
    or 'such that', which starts
    the list of constraints.
    \item $b_j^{\text{L}}$ and $b_j^{\text{U}}$ are called 
    the lower bound and upper bound on
    constraint $j$, respectively.
    \item If $b_j^{\text{L}} = b_j^{\text{U}}$, the constraint is said to be an equality constraint, and
    can be represented using $=$ sign.
    \item A constraint can be represented with only one bound. In this case,
    the other bound is either $b_j^{\text{L}} = -\infty$ or 
    $b_j^{\text{U}} = +\infty$
\end{itemize}

Definitions for objectives:
\begin{itemize}
    \item The function (1) is the \emph{objective function}.
    \item $\zeta$ represents the value of the objective 
    function, called \emph{objective value}. For a given 
    solution $x = (x_1, x_2, \dots, x_n)$, we use $\zeta(x)$ to
    denote the objective value of $x$:
    $\zeta(x) = c_1x_1 + c_2x_2 + \dots + c_nx_n + c_0$
    \item The symbol $\text{min}$ or $\text{max}$ indicates 
    that we want to minimize or maximize the objective value $\zeta$. 
    \item A solution $x^*$ is \emph{optimal} if it is feasible and 
    $\zeta(x^*) \leq \zeta(x')$ for any other feasible solution $x'$
    (for minimization).
    \item The objective value of the optimal solution is called 
    \emph{optimal value} and is denoted by $\zeta(x^*)$ or $\zeta^*$.
\end{itemize}

An LP has three possible outcomes:
\begin{itemize}
    \item LP is \emph{optimal}, which means that there exists an optimal solution,
    either uniquely or among infinitely many others.
    \item LP is \emph{infeasible}, which means that a feasible solution does not
    exist.
    \item LP is \emph{unbounded}, which means that the objective value goes to
    infinity (for maximization). In particular, for any real number $K$,
    there always exists a feasible solution $x$ such that $\zeta(x) \geq K$.
\end{itemize}

LPs can be represented in matrix form. The 
general form corresponding to Equation \ref{eq:genlp}
is 
\[
\begin{aligned}
\min_{x}\quad & \zeta = c^\top x \\
\text{s.t.}\quad & b^L \le Ax \le b^U, \\
& l \le x \le u.
\end{aligned}
\]

Here,
\[
c =
\begin{bmatrix}
c_1\\
\vdots\\
c_n
\end{bmatrix},
\qquad
x =
\begin{bmatrix}
x_1\\
\vdots\\
x_n
\end{bmatrix},
\qquad
A =
\begin{bmatrix}
a_{1,1} & \cdots & a_{1,n}\\
\vdots & \ddots & \vdots\\
a_{m,1} & \cdots & a_{m,n}
\end{bmatrix},
\]
\[
b^L =
\begin{bmatrix}
b^L_1\\
\vdots\\
b^L_m
\end{bmatrix},
\qquad
b^U =
\begin{bmatrix}
b^U_1\\
\vdots\\
b^U_m
\end{bmatrix},
\qquad
l =
\begin{bmatrix}
l_1\\
\vdots\\
l_m
\end{bmatrix},
\qquad
u =
\begin{bmatrix}
u_1\\
\vdots\\
u_m
\end{bmatrix}.
\]